% Volume~III, appendix C: logs (English version). Translator EN2, 2026-10-29, from
% memoire/tome3/fr/annexes/C-journaux.tex (writer D); same labels, same macros, same numbers.
% Model: the appendix on the logs of Volume I (memoire/annexes/journaux.tex). The table of folders is generated by
% memoire/outils/tome3_annexes.py from commit 4b071d0 (memoire/tome3/annexes/journaux-tome3.tex, bilingual by \EF);
% the sessions and the cut passes come from the deliverables, the judgements and the resumption notes cited in
% their rows.

\chapter{Logs}\label{t3ann:journaux}

\begin{origine}
\textbf{At the start}, the rule of the first two volumes: whatever produces a published number is committed, tools and logs.
\textbf{This appendix establishes} where the logs of Volume~III are, folder by folder; the sessions that carry a speed measurement, with their conditions; and the passes cut, discarded or published as inconclusive, with their reason, because none disappears silently.
\textbf{What remains} are the heavy files outside the repository (f64 references, \texttt{.kld}), which can be regenerated by the committed tools.
\end{origine}

\section{The log folders}\label{t3ann:sec:dossiers}

\Cref{t3ann:tab:journaux} gives each folder of logs, CSVs and notes committed for this volume, its number of files and its size (the tools are in appendix~\ref{t3ann:outils}). It is generated by \fichier{memoire/outils/tome3_annexes.py} from the commit: the size of a folder is the sum of the sizes of the git objects of its files (\texttt{git ls-tree -l}), independent of the line endings of the copy (corrected on 2026-10-29, appendix~\ref{t3ann:outils}). The folders of the leads (\fichier{mesures/pistes/}) and of mimicry belong to their streams; those of the disk placement judge (\fichier{mesures/disque/placement/juge/}) are there for reserve R1. Those of the second draw of P3, in progress elsewhere, are not in it (R6).

\input{\Racine/tome3/annexes/journaux-tome3}

\paragraph{What is not in the repository.} The heavy files: the f64 references and the \texttt{.kld} of the KL passes (about \SI{13}{\Go} for the judge of the construction alone, under \fichier{~/dev/essais/troncons/}), the \texttt{.sym} files of the 768 experts and the fast companions. Each is regenerated by a committed tool, and its fingerprint is in a committed log when a figure depends on it (for example the output \texttt{.kld} of J-C1, \cref{t3jug:sec:jc1}).

\section{The sessions that carry a speed figure}\label{t3ann:sec:seances}

A speed figure is worth only as much as its session: the machine, the power supply, the load of the other streams, the card lock. \Cref{t3ann:tab:seances} gives those of the volume.

\begin{table}[ht]
\centering\footnotesize
\caption[The speed sessions of Volume~III]{The sessions that carry a speed figure of Volume~III, their conditions and what they decide \mes.}
\label{t3ann:tab:seances}
\begin{tabularx}{\linewidth}{@{}>{\raggedright\arraybackslash}p{0.17\linewidth}>{\raggedright\arraybackslash}p{0.20\linewidth}>{\raggedright\arraybackslash}X>{\raggedright\arraybackslash}p{0.22\linewidth}@{}}
\toprule
Session & When & Conditions & What it decides \\
\midrule
lot B, s1 (20b) & 17 Oct., start at 11:45 & A B B A A B B A; under the load of the KL pass B0, of the computation of the rate and of the writing of the companion of the 120b ($\sigma_A = \SI{17.8}{\percent}$) & published, inconclusive \\
lot B, s2 (20b) & 22 Oct., 1:11 -- 1:25 & A B B A A B B A; two empty readings at the start, no pass under load, mains power, lock ($\sigma_A = \SI{2.3}{\percent}$) & $\times 1.24$ announced, struck through as a gain of the mode (G2) \\
lot B, 120b-V & 22 Oct., 2:34 -- 3:38 & same rules ($\sigma_A = \SI{6.6}{\percent}$) & $\times 2.00$ announced, struck through as a gain of the mode (G1) \\
judge, three-file session (120b then 20b) & 22 Oct., engine from 6:17 to 7:11 & A B R R B A A B R R B A, disk stage first; none of the 36 passes under load; mains power at each pass; lock held & $r(R/B) = \num{1.284}$ and \num{0.942}; J1 given again (\num{1.442}) \\
\bottomrule
\end{tabularx}
\sourcefig{\fichier{mesures/genome/construction/synonyme/vitesse/pilote-s1.json}, \fichier{pilote-s2.json}, \fichier{pilote-s2-120b.json}; \fichier{mesures/genome/construction/juge/vitesse/seance-\{120b,20b\}.\{json,log\}}. The session of lot A is in the chapter on the codon; those of the disk placement, in the chapter ``The disk stage'' planned for Volume~II.}
\end{table}

\section{The passes cut, discarded, or inconclusive}\label{t3ann:sec:coupees}

No pass disappears silently: those that decided nothing are published as such, with their reason.
\begin{itemize}
\item \textbf{Session s1 of lot B}: published, inconclusive, because it ran under the load of other passes; redone on a calm machine (s2), and this is what led to the recipe ``a speed session starts only on a machine recorded as calm'' (\cref{t3sno:sec:seances}).
\item \textbf{The first B pass of session s2} (\num{5.46} tokens/s, just after the companion was written): kept in the median, as the declared rule required; the judge shows that with the noise of the B passes it would have prevented the announcement (finding M4).
\item \textbf{The judge's substituted f64 pass $\lambda = 0$}: the bash loop that launched it was stopped, but the native program finished on its own; its redirected output was complete and was kept under the name \fichier{sub-l0.json}. The substituted f64 pass $\lambda = 1$ was not launched (time, and the placement stream waiting for calm), and this is stated at each verdict that depends on it (\cref{t3jug:sec:klnombres}).
\item \textbf{The broken copies} (the flipped bias of J-C1, the breaks of the figures) are erased after use; their logs remain.
\item \textbf{The second draw of P3}, in progress in another session, has attempts cut by external causes (mains power, the machine's memory, disk), without any usable measurement and without any threshold changed; its logs are committed by its session, and the volume draws no figure from them (reserve R6).
\end{itemize}

\section*{What could be wrong}

\begin{faux}
\begin{itemize}
\item \textbf{The size of a folder} is that of the committed files, with LF line endings (a CRLF copy is larger); it says nothing of the heavy regenerable files outside the repository.
\item \textbf{The sorting by folder} follows the depth of the paths (as in Volume~II); a folder that mixes two steps appears in it as a single row.
\item \textbf{``No pass under load''} is what the declared reading says (more than \SI{5}{\percent} of a core or more than \SI{8}{\Mo} read in \SI{3}{\second}); a smaller load, or one between two readings, escapes it.
\end{itemize}
\end{faux}
